\uuid{pPkI}
\titre{Nombres premiers}
\auteur{bodin ; croizat ; sacre}
\organisation{exo7}
\licence{CC BY-NC-SA 4.0}
\datecreate{2022-03-08}
\niveau{L1}
\module{Algèbre}
\chapitre{Arithmétique dans Z}
\sousChapitre{Nombres premiers, nombres premiers entre eux}
\difficulte{2}
\theme{Nombres premiers}
\exo7qcmid{hilisit.arithmetique.033}

\contenu{
\begin{question}
\qclassification{103.04}
    À l'aide d'une calculatrice, quelle est l'écriture de la décomposition en produit de facteurs premiers de $N = 111 \, 111$ ?
    \begin{answers}
        \bad{$N = 11 \times 10\,101$.}
        \bad{$N = 3 \times 11 \times 3367$.}
        \bad{$N = 7 \times 33 \times 481$.}
        \good{$N = 3 \times 7 \times 11 \times 13 \times3713$.}
    \end{answers}
    \begin{explanations}
    Les entiers $10\,101$, $3367$ et $481$ sont des multiples de $13$ !
    \end{explanations}
\end{question}
}
